\documentclass[11pt]{article}
\usepackage[a4paper,margin=1in]{geometry}
\usepackage{amsmath,amssymb,mathtools,bm}
\usepackage{enumitem,booktabs,array}
\usepackage{tcolorbox}
\usepackage{hyperref}
\usepackage[protrusion=true,expansion=false]{microtype}
\hypersetup{colorlinks=true,linkcolor=blue,urlcolor=blue,citecolor=blue}
\setlist[itemize]{topsep=2pt,itemsep=2pt}
\setlist[enumerate]{topsep=2pt,itemsep=2pt}
\newtcolorbox{keybox}{colback=blue!3!white,colframe=blue!40!black,boxrule=0.4pt,arc=1mm}
\newtcolorbox{alertbox}{colback=red!3!white,colframe=red!40!black,boxrule=0.4pt,arc=1mm}
\newtcolorbox{answerbox}{colback=green!3!white,colframe=green!40!black,boxrule=0.4pt,arc=1mm}
\newtcolorbox{drillbox}{colback=yellow!5!white,colframe=orange!60!black,boxrule=0.4pt,arc=1mm}
\newcommand{\EV}{\mathbb{E}}
\newcommand{\Var}{\mathrm{Var}}
\newcommand{\Cov}{\mathrm{Cov}}
\newcommand{\blank}[1]{\underline{\hspace{#1}}}

\title{E02-03: Azar, Schmalz \& Tecu (2018, JF) \\
\large ``Anticompetitive Effects of Common Ownership'' --- Active Recall Workbook}
\author{Empirical Corporate Finance II --- Exam Prep}
\date{\today}
\begin{document}\maketitle

\begin{keybox}
\textbf{Paper in one sentence.} Using US airline route-level data 2001--2014, AST show that ticket prices are 3--7\% higher on routes where competing airlines share more institutional-investor owners, and propose a Modified-HHI (MHHI) that measures market power under common ownership --- suggesting that diversified index funds create implicit anticompetitive incentives.

\textbf{Placement.} Canonical ``common ownership'' paper in corporate finance and IO. Generated the antitrust debate on horizontal shareholding. Paired with Lewellen-Lowry (E02-04) which challenges the coordination mechanism.
\end{keybox}

\section*{Round 1: Complete Walk-Through}

\subsection*{1.1 Setup}
\begin{itemize}
\item Route level panel of US domestic airlines, 2001--2014, quarterly.
\item Institutional-holdings data (13F) identify which institutions hold each airline.
\item Construct the Modified Herfindahl-Hirschman Index (MHHI) that adds a ``common-ownership'' term to standard HHI.
\end{itemize}

\subsection*{1.2 Modified HHI}
Let $s_i$ be firm $i$'s market share and $\beta_{if}$ institution $f$'s equity stake in firm $i$. Then
\[
\text{MHHI} = \underbrace{\sum_i s_i^2}_{\text{HHI}} + \underbrace{\sum_{i\ne j} s_i s_j \,\frac{\sum_f \beta_{if}\beta_{jf}}{\sum_f \beta_{if}^2}}_{\text{common ownership term}}.
\]
The second term captures the degree to which rival firms $i,j$ share owners.

\subsection*{1.3 Identification}
\begin{itemize}
\item Panel with route + time FE.
\item Primary worry: MHHI correlates with unobserved demand. Instruments: BlackRock-BGI merger (2009) as exogenous shock to common ownership.
\item IV sign supports causal effect of MHHI on prices.
\end{itemize}

\subsection*{1.4 Headline results}
\begin{itemize}
\item Moving from zero common ownership to sample average raises route prices 3--7\%.
\item Effect robust to route FE, time FE, standard covariates.
\item IV using BlackRock-BGI merger: airlines with greater post-merger common ownership raise prices more on affected routes.
\end{itemize}

\subsection*{1.5 Proposed mechanisms}
\begin{itemize}
\item Managers internalize common owners' preference for higher industry profits (not firm-specific profits).
\item Through voting, engagement, and compensation committees, common owners tilt firms toward softer competition.
\item Informal norms, analyst conferences, and industry meetings also transmit preferences.
\end{itemize}

\subsection*{1.6 Caveats}
\begin{alertbox}
\begin{itemize}
\item Mechanism not directly observed; correlation vs.\ causation debated (see Lewellen-Lowry 2021).
\item MHHI construction sensitive to normalization choices.
\item Price effects may reflect alternative channels (quality, network structure).
\end{itemize}
\end{alertbox}

\section*{Round 2: Level 1 Fill-in}
\begin{enumerate}
\item Sample: US \blank{1.5cm} routes, 2001--\blank{1cm}.
\item Index constructed: \blank{1cm}HHI.
\item Price effect from zero to mean common ownership: \blank{1cm}--\blank{1cm}\%.
\item Key IV: \blank{1.5cm}--BGI merger 2009.
\item Proposed mechanism: managers internalize \blank{2cm} owners' preferences.
\end{enumerate}

\section*{Round 3: Level 2 Prompts}
\begin{enumerate}
\item Define MHHI and state the common-ownership term's intuition.
\item Describe the empirical design and identification strategy.
\item Discuss proposed mechanisms for how common ownership affects prices.
\item Relate AST to Lewellen-Lowry (2021).
\item Evaluate antitrust-policy implications.
\end{enumerate}

\section*{Round 4: Minimal Scaffolding}
From a blank page: (i) setting, (ii) MHHI formula, (iii) identification, (iv) results, (v) mechanism.

\section*{Round 5: Blank Page}
Essay: common ownership and product-market competition --- evidence from airline routes.

\vspace{6pt}
\begin{drillbox}
\textbf{Speed Drill.} (1) Data. (2) MHHI formula. (3) Price effect. (4) IV. (5) Mechanism claim.
\end{drillbox}

\section*{Quick Reference \& Answer Key}
\begin{answerbox}
\begin{itemize}
\item \textbf{Setting.} US airline routes, 2001--14.
\item \textbf{Tool.} MHHI adds common-ownership term.
\item \textbf{Effect.} $+3$--$7$\% prices.
\item \textbf{IV.} BlackRock-BGI merger.
\end{itemize}
\end{answerbox}

\begin{keybox}
\textbf{30-second exam pitch.} Azar, Schmalz \& Tecu (2018) document that US airline-route ticket prices are 3--7\% higher where competing airlines share more institutional-investor owners. They construct a Modified Herfindahl-Hirschman Index (MHHI) that augments standard market-concentration measures with a common-ownership term reflecting how much rivals share owners. Using airline route-quarter panels 2001--2014 with route and time fixed effects, and an IV from the 2009 BlackRock-BGI merger, they argue the relationship is causal: more common ownership raises ticket prices. Proposed mechanism: managers internalize common owners' preference for softer industry-wide competition through voting, engagement, and compensation. The paper launches the 2010s common-ownership debate in antitrust and corporate governance, motivating policy proposals to restrict horizontal shareholding. It is paired with Lewellen-Lowry (2021), which questions whether common owners actively coordinate firms.
\end{keybox}

\end{document}
